\section{The algebra and estimates of an ordinary sweep}\label{sec:sweep}

An ordinary sweep inserts a potential increment at each active site while
decreasing the artificial damping. The physical sites remain active,
and the vector $r$ remains frozen, so we must keep the endpoint
diagonal error small.
The random potential increment is
conditionally centered; its nonlinear effect on the resolvent has a small
conditional mean, but an individual update need not be small in every useful norm.
We therefore estimate a whole family of possible next updates at a fixed
prestate. Random ordering will turn these estimates into bounds for one
tile in Section~\ref{sec:local-sweeps}.

We first explain the change of coordinates that expresses a bit-port deletion
as a Schur complement plus an explicit correction. We then control three quantities: the matrix
$K$, the Laplacian $D$, and the diagonal self-consistency error. A fourth
estimate preserves the row fourth powers needed in all three arguments.
Throughout Section~\ref{sec:sweep}, $d\ge3$ is fixed, $N=L^d$ with
$L=4^{20m}$, and constants may depend on $d$ and the broad bounds in
Section~\ref{sec:state}. In particular, they do not depend on the number
of ports or on the subset of possible next updates.

All the one-step identities and estimates used below are proved here for
the present $d$-dimensional operator. We use the parameters
$\delta=.01$, $\kappa=.005\delta$, and $\nu=10^{-5}\delta/d$ of
Section~\ref{sec:state}; $\nu$ enters the later geometry, not the
one-step algebra.

Self-consistent resolvent equations and fourth-power row estimates as
inputs to concentration also appear in
\cite[Sections~1.1--1.2]{BlackDroginHernandez2025}.
Here we give explicit one-step estimates for every subset of remaining
ports.

\subsection{Splitting a site into two ports}

At the beginning of a stage $n<m$ that inserts one digit, put $t=t_n$ and
\[
 w=\frac{3s_n}{4t_n},\qquad
 \theta=1-w=\frac{t_{n+1}}{t_n}.
\]
For the final residual sweep at $n=m$, take instead $t=t_m=2\eta$
and $w=\theta=1/2$. In both cases $\theta\in[1/4,5/8]$.
The one-step estimates below use only this range and the conditional
moment assumptions stated below, so they apply to the final suffix as
well as to a single digit.

Start from a committed physical state at a stage boundary satisfying the
bounds of Section~\ref{sec:state}, including $r_i^{\rm phys}=g_{ii}^{\rm phys}$.
The terminal set stays fixed throughout the sweep. At a numerical prefix,
let $M_{\rm cur}$ be the real Schur matrix with the current real potential
values, and let $\mathcal I$ be the set of active physical
sites whose current sweep increment has already been inserted. The physical
inverse represented at that prefix is
\begin{equation}\label{eq:sweep-physical-prefix}
 g^{\rm phys}
 =\left(M_{\rm cur}-t\diag\bigl((\vartheta_i r_i^{\rm phys})_{i\in A}\bigr)\right)^{-1},
 \qquad
 \vartheta_i=\theta+w\mathbf1_{\{i\notin\mathcal I\}}.
\end{equation}
Thus an unprocessed site retains the full frozen shift $-tr_i^{\rm phys}$,
whereas a processed site retains $-\theta t r_i^{\rm phys}$. The matrix
$M_{\rm cur}$ records the real potential additions; $r^{\rm phys}$ does not
change. This description applies equally to a committed prefix and to the
current provisional prefix of a tile trial.

Represent each active physical site initially by a bit port and a reserve
port, with weights $k_{\rm bit}=w^{1/4}$ and
$k_{\rm res}=\theta^{1/4}$. The bit port is deleted when that site is
processed, and the reserve port remains throughout the sweep. In particular,
deleting a bit port does not make its physical site terminal. For the ports
remaining above $i$, one has $\sum_x k_x^4=\vartheta_i$.
If $x$ is a port above $i$ and $y$ is a port above $j$, define
\begin{equation}\label{eq:sweep-virtual}
 \begin{aligned}
 g_{xy}&=k_xk_y g^{\rm phys}_{ij},&
 r_x&=k_x^2r^{\rm phys}_i,& p_x&=\Im r_x,\\
 h&=\sqrt t\,g,& K_{xy}&=h_{xy}^2,& F_{xy}&=|h_{xy}|^2.
 \end{aligned}
\end{equation}
After a deletion these formulas use the new physical inverse and the
remaining ports. This virtual $g$ is not an inverse on the duplicated
space; its entries are repeated and rescaled physical resolvent entries.
Put $q_x=\Im g_{xx}$ and
$D=\diag(pq)-\diag(p)F\diag(p)$. When $D$ has only the constant kernel,
$X=D^+$ denotes its inverse form on neutral port charges, as in
Section~\ref{sec:state}.

\begin{lemma}[Port identities]\label{lem:duplication}
For every numerical prefix represented by
\eqref{eq:sweep-physical-prefix}--\eqref{eq:sweep-virtual}, the virtual
matrices satisfy the Ward identity
$h\diag(p)h^*=\Im g$. Their operator norm obeys
$\|h\|\le C/\sqrt t$. At the initial duplication, the operator norm
of $K$ is unchanged. The boundary fourth-power bound implies that
every full fourth-power row sum is at most $3H_*/4$, and every reserve
row summed only over reserve columns is at most $\theta^2H_*$.

For a potential $f$ on the initial ports, set
$\bar f_i=\sum_{x\text{ above }i}k_x^4 f_x$. Then
\begin{align}
 D_0(f)
 &=D^{\rm phys}(\bar f)
   +\sum_i(p_i^{\rm phys})^2
       \sum_{x\text{ above }i}k_x^4|f_x-\bar f_i|^2,
       \label{eq:sweep-duplicate-energy}\\
 X_0(y)
 &=X^{\rm phys}(y_\Sigma)
   +\sum_x\frac{|y_x|^2}{k_x^4(p_i^{\rm phys})^2}
   -\sum_i\frac{|(y_\Sigma)_i|^2}{(p_i^{\rm phys})^2},
       \label{eq:sweep-duplicate-dual}
\end{align}
where the second identity is for neutral charges and
$(y_\Sigma)_i=\sum_{x\text{ above }i}y_x$.
In particular $X_0\asymp W^*$ and $D_0\asymp W$, and the sharp
stage-start comparison implies
\begin{equation}\label{eq:sweep-sharp-duplicate}
 X_0\le (1+t^\kappa)\diag(p_x^{-2})+C_n t\Lambda^+_\Sigma.
\end{equation}
\end{lemma}

\begin{proof}
Let $V_{xi}=k_x\mathbf1_{\{x\text{ above }i\}}$ on the remaining
ports. The physical imaginary damping at $i$ is
$t p_i^{\rm phys}\sum_{x\text{ above }i}k_x^4\ge ct$.
Moreover,
\[
 V^*\diag(p_x)V
 =\diag\left(p_i^{\rm phys}
                  \sum_{x\text{ above }i}k_x^4\right).
\]
Applying the physical Ward identity to $g=Vg^{\rm phys}V^\top$
proves the virtual one. Since $p\ge c$,
$c\|h\|^2\le\|\Im h\|/\sqrt t\le\|h\|/\sqrt t$, proving the norm
bound.

At the initial duplication, the matrix
$Q_{xi}=k_x^2\mathbf1_{\{x\text{ above }i\}}$ satisfies $Q^*Q=\Id$,
and $K_0=QK^{\rm phys}Q^\top$. Also
\[
 \sum_y|h_{xy}|^4
 =k_x^4\sum_j|h^{\rm phys}_{ij}|^4,
\]
because the fourth powers of the two weights over a site sum to one.
The stated bounds follow from $\max(w,\theta)\le3/4$; restricting the
columns to reserves introduces a further factor $\theta$.

At a stage start, $q_i^{\rm phys}=p_i^{\rm phys}$. Expanding the two
terms in $D_0$ gives \eqref{eq:sweep-duplicate-energy}. Its dual
formula follows by maximizing
$2\Re\langle y,f\rangle-D_0(f)$ separately over the averages
$\bar f$ and the weighted-mean-zero fluctuations above each site.
The broad comparison with $W^*$, and its dual with $W$, follow from
the physical comparisons and the fact that the two weights are bounded
away from zero.

For the sharp bound, retain the exact diagonal cancellation in
\eqref{eq:sweep-duplicate-dual}. Weighted Cauchy--Schwarz gives
\[
 \frac{|(y_\Sigma)_i|^2}{(p_i^{\rm phys})^2}
 \le\sum_{x\text{ above }i}
       \frac{|y_x|^2}{k_x^4(p_i^{\rm phys})^2}.
\]
The coefficient of the diagonal form is therefore exactly
$1+t^\kappa$, as claimed.
\end{proof}

The reserve part of this identity is the target of the sweep. If $y$ is a
neutral physical charge, let $y_{\rm res}$ put that charge on the reserve
ports and zero on the bit ports. The exact initial identity is
\begin{equation}\label{eq:sweep-reserve-target}
 \theta X_0(y_{\rm res})
 =\theta X^{\rm phys}(y)
   +(1-\theta)\sum_i\frac{|y_i|^2}{(p_i^{\rm phys})^2}.
\end{equation}
If all bit ports have been deleted, the definitions give
$h_{\rm cur}=\sqrt{\theta t}\,g^{\rm phys,new}$ and
$p_{\rm cur}=\sqrt\theta\,p^{\rm phys}$. When the current Laplacian has
only the constant kernel, these scalings give
$X^{\rm phys,new}(y)=\theta X_{\rm cur}(y_{\rm res})$.
We must therefore keep $X_{\rm cur}$ close to $X_0$ on the surviving
reserves. In the target form \eqref{eq:sweep-reserve-target}, an old sharp
diagonal coefficient $1+t^\kappa$ becomes $1+\theta t^\kappa$.
Since $\theta^\kappa>\theta$, this lies below the next coefficient
$1+(\theta t)^\kappa$, leaving room for later errors. The initial
reserve-row fourth-power sum,
$\theta^2H_*$, likewise has strict slack below the cap at the new scale,
$(\theta t)^{1+\delta}$. The estimates below control the corrections
that could consume these two margins, as well as the $K$ norm and the
diagonal error.

\subsection{One deletion and the transported matrices}

Let $\mathscr F$ denote the retained history just before choosing the
next permutation slot, including the numerical prestate and all earlier
queries. Conditional on choosing a particular remaining bit port $x$
above a physical site $i$, the real variable $\xi_0$ is required to satisfy
\begin{equation}\label{eq:sweep-digit-moments}
 \EE(\xi_0\mid\mathscr F,x)=0,\qquad
 \EE(\xi_0^2\mid\mathscr F,x)=t,\qquad |\xi_0|\le C\sqrt t.
\end{equation}
Deleting $x$ adds the real potential increment $k_x^2\xi_0=\sqrt w\,\xi_0$ to the
partial potential at $i$. It also changes $\vartheta_i$ from $1$ to
$\theta$ in \eqref{eq:sweep-physical-prefix}, thereby adding
$tw r_i^{\rm phys}=k_x^2tr_x$ to the diagonal of the damped matrix.
The complete change to that matrix is therefore
$k_x^2(\xi_0+tr_x)$. Only the first summand changes the real partial
potential; the second removes the bit port's share of the frozen shift
and damping.
For an ordinary stage, $\xi_0$ is the next true digit, or an auxiliary
digit when prescribed, divided by $\sqrt w$; its physical variance is
$s_n-s_{n+1}=wt$. The expectation symbol in the remainder of
Section~\ref{sec:sweep} denotes this conditional expectation whenever
the prestate and hypothetical pivot have been fixed. In such statements,
``conditioning on the prestate'' always includes the full retained
history $\mathscr F$, not just the matrix entries of the current state.
At the final sweep, the untouched true suffix $\zeta$ has conditional
mean zero, variance $s_m=\eta=wt$, and $|\zeta|\le\sqrt{3\eta}$;
thus $\xi_0=\zeta/\sqrt w$ satisfies
\eqref{eq:sweep-digit-moments}. We assume neither symmetry nor two-point
support in any estimate below. The ordinary query rule is verified in
Section~\ref{sec:local-sweeps} using Lemma~\ref{lem:conditional-freshness};
the conditional law is an input here, and no conclusion conditions on
future success of the construction.
The Laplacian pivot is strictly positive: by Ward,
$q_x-p_xF_{xx}=\sum_{y\ne x}p_yF_{xy}$, and the reserve port at the
same site contributes a positive term because the physical diagonal
inverse has positive imaginary part.

\begin{lemma}[Single deletion]\label{lem:single-update}
For sufficiently small $t$, suppose $|g_{xx}|+|r_x|\le C$ and
$q_x-p_xF_{xx}>0$, and let $R$ be the remaining ports
other than $x$. With $l=h_{Rx}$,
\begin{equation}\label{eq:sweep-single-update}
 h'=h_{RR}-\gamma ll^\top,\qquad
 \gamma=\frac{\xi_0+tr_x}{\sqrt t[1+(\xi_0+tr_x)g_{xx}]}.
\end{equation}
Put $z_x=\sqrt t(g_{xx}-r_x)$ and $a_x=1-K_{xx}$. Then
\begin{equation}\label{eq:sweep-gamma-moments}
 |\gamma|\le C,\quad
 \EE\gamma=-z_x+O(t),\quad
 \EE\gamma^2=1+O(\sqrt t),\quad
 \EE|\gamma|^2=1+O(\sqrt t).
\end{equation}
On the survivor coordinates,
\begin{align}
 \Id-K'&=\Schur_{\{x\}}(\Id-K)-J_x,\label{eq:sweep-j-identity}\\
 D'&=\Schur_{\{x\}}(D)+L_x,\label{eq:sweep-l-identity}\\
 J_x&=-2\gamma\diag(l)h_{RR}\diag(l)
          +(\gamma^2-a_x^{-1})l^2(l^2)^\top.
          \label{eq:sweep-j-formula}
\end{align}
Here $l^2$ is the entrywise square. The real symmetric Laplacian
$L_x$ has survivor edge weights
\begin{equation}\label{eq:sweep-l-formula}
 p_ip_j\left[
 -2\Re(\gamma\bar h_{ij}h_{ix}h_{xj})
 +\left(|\gamma|^2-\frac{p_x}{q_x-p_xF_{xx}}\right)F_{ix}F_{xj}
 \right].
\end{equation}
The subscript of $\Schur$ denotes the indices eliminated.
\end{lemma}

\begin{proof}
Sherman--Morrison on the physical inverse, followed by
\eqref{eq:sweep-virtual}, gives \eqref{eq:sweep-single-update}.
For the moment calculation write $\varepsilon=\sqrt t$ and
$\xi_0=\varepsilon\xi$. The denominator is $1+O(\varepsilon)$ uniformly,
and
\[
 \gamma=\xi+\varepsilon(r_x-g_{xx}\xi^2)
       +\varepsilon^2(g_{xx}^2\xi^3-2r_xg_{xx}\xi)
       +O(\varepsilon^3).
\]
Now use $\EE\xi=0$, $\EE\xi^2=1$ and boundedness of $\xi$.
No vanishing third moment is needed.

Entrywise squaring \eqref{eq:sweep-single-update} and using the
scalar pivot $a_x$ in $\Id-K$ gives
\eqref{eq:sweep-j-identity}--\eqref{eq:sweep-j-formula}.
For $D$, the old edge conductance is $p_ip_jF_{ij}$ and its pivot
is $p_x(q_x-p_xF_{xx})$. Eliminating $x$ adds conductance
$p_ip_jp_xF_{ix}F_{xj}/(q_x-p_xF_{xx})$.
Subtracting this from the conductance obtained by expanding
$|h_{ij}-\gamma h_{ix}h_{xj}|^2$ proves
\eqref{eq:sweep-l-formula}.
\end{proof}

Pad $J_x$ and $L_x$ by zero on the deleted coordinate and on all
earlier deleted ports. Adding such a survivor-supported matrix commutes
with earlier elimination. Writing $E_{\rm del}$ for the set of deleted
ports, induction gives the exact cumulative identities
\begin{equation}\label{eq:sweep-cumulative-schur}
 \Id-K_{\rm cur}=\Schur_{E_{\rm del}}\left(\Id-K_0-\sum J_x\right),\qquad
 D_{\rm cur}=\Schur_{E_{\rm del}}\left(D_0+\sum L_x\right).
\end{equation}
For later use impose the prefix bounds
\begin{equation}\label{eq:sweep-global-caps}
 \left\|\sum J_x\right\|\le t^{.10\delta},\qquad
 \left|\sum L_x\right|\le t^{.10\delta}D_0.
\end{equation}
The second bound is a two-sided quadratic-form bound.
These caps imply $D_{\rm cur}(f|_{\rm cur})\le C W(f)$ for every
extension $f$ to the initial ports, and
$X_{\rm cur}=(1+O(t^{.10\delta}))X_0$ on survivor-supported neutral
charges.

They also control $K$, despite the fact that $\Id-K$ is complex
symmetric rather than self-adjoint. Indeed, let $A=\Id-B$ with
$\|B\|\le\rho<1$, and write
$\Id-\widehat B=\Schur_{E_{\rm del}}A$ on the survivors.
The harmonic extension of a survivor vector $v$ to $u$ exists because
$\|B\|<1$; it satisfies $u_{\rm cur}=v$ and
$(Au)_{\rm deleted}=0$. Hence $(Bu)_{\rm deleted}=u_{\rm deleted}$
and $(Bu)_{\rm cur}=\widehat Bv$, so
\[
 \|u_{\rm deleted}\|^2+\|\widehat Bv\|^2
 \le\rho^2(\|u_{\rm deleted}\|^2+\|v\|^2).
\]
Consequently $\|\widehat B\|\le\rho$. Apply this with
$B=K_0+\sum J_x$ and $\widehat B=K_{\rm cur}$.

\subsection{Consequences of the prestate caps}

During a sweep set
\[
 \beta=.42\delta,\qquad
 P=(\Id-K)^{-1},\qquad
 z=\sqrt t(g_{\rm diag}-r),\qquad
 U=Pz,\qquad \widetilde U=U/\sqrt t.
\]
We call $g_{\rm diag}-r$ the raw diagonal error and
$\widetilde U=P(g_{\rm diag}-r)$ its transformed lag. On the current
ports their exact relation is
\[
 g_{\rm diag}-r=(\Id-K)\widetilde U.
\]
On the virtual space we use this algebraic identity. Comparing the equation
$(\Id-K)U=z$ before and after a deletion
will combine the diagonal change with the Schur change of $\Id-K$.
Lemma~\ref{lem:sweep-lag} uses the cancellation
$\EE\gamma=-z_x+O(t)$ in that combined equation to control the
transformed lag. At the start $U=0$ by the physical fixed-point relation.

In addition to \eqref{eq:sweep-global-caps},
the prestate caps are
\begin{equation}\label{eq:sweep-row-lag-caps}
 \max_i\sum_jF_{ij}^2\le H_*=t^{1+\delta},
 \qquad \|\widetilde U\|_\infty\le t^\beta.
\end{equation}
The frozen $r$ has bounded modulus and its imaginary part lies between
fixed positive constants. There are at least $N/2$ active physical
sites. The local stopping rule will retain an update only when the two
displays of mutable caps hold at its prestate; the stage induction supplies
the fixed bounds on $r,p$ and the active-site count.

\begin{lemma}[Analytic prestate bounds]\label{lem:sweep-prestate}
Suppose the initial physical state satisfies the bounds at a stage boundary in
Section~\ref{sec:state}, including self-consistency, and is duplicated as
in Lemma~\ref{lem:duplication}. Suppose its current numerical prefix satisfies
\eqref{eq:sweep-global-caps}--\eqref{eq:sweep-row-lag-caps}, the broad
bounds on the frozen $r,p$, and the active-site count above. Then
\begin{gather}
 |g_{ii}-r_i|\le Ct^\beta,\quad q_i=p_i+O(t^\beta),\quad
 \sum_jF_{ij}\le C,\quad |h_{ii}|\le C\sqrt t,\label{eq:sweep-broad}\\
 T_*^2:=\max_i\|\diag(h_i)h\|^2
       \le C\sqrt{H_*/t}=Ct^{\delta/2},\label{eq:sweep-t-bound}\\
 (F^3)_{ii}\le Ct^\kappa H_*.\label{eq:sweep-f-cube}
\end{gather}
The bounds hold uniformly for every current set of remaining ports.
\end{lemma}

\begin{proof}
Since $z=(\Id-K)U$ and $Fp=q$,
\[
 |g_{ii}-r_i|
 \le t^\beta\left(1+\frac{q_i}{\min_jp_j}\right).
\]
Taking imaginary parts first bounds $q_i$ for small $t$; substituting
back proves \eqref{eq:sweep-broad}.

For \eqref{eq:sweep-t-bound}, Ward gives
\[
 \diag(h_i)hh^*\diag(\bar h_i)
 \le Ct^{-1/2}\diag(h_i)(\Im h)\diag(\bar h_i).
\]
Apply the weighted Schur test to the matrix on the right using weights
$|h_{ij}|$, omitting its zero rows and columns. Its weighted absolute
row quotient is at most
\[
 \sum_k|h_{jk}|F_{ik}
 \le\left(\sum_kF_{jk}\right)^{1/2}
       \left(\sum_kF_{ik}^2\right)^{1/2}
 \le C\sqrt{H_*}.
\]
This proves the operator estimate.

To prove \eqref{eq:sweep-f-cube}, put $M=\diag(pq)$, fix $i$, and
define
\[
 f_j=\frac{F_{ij}}{p_j},\qquad
 \bar m=\frac{\sum_jM_{jj}f_j}{\Tr M},\qquad
 y=M(f-\bar m\mathbf1).
\]
The charge $y$ is neutral. The variational formula for $X$ gives
\[
 X(y)\ge2\langle y,f\rangle-D(f)
       =f^\top Mf-2\bar m^2\Tr M+(F^3)_{ii}.
\]
On the other hand, \eqref{eq:sweep-sharp-duplicate}, Schur duality,
and the lower half of \eqref{eq:sweep-global-caps} imply
\[
 X\le(1+Ct^\kappa)M^{-1}+Ct\Lambda^+_\Sigma.
\]
Here we used $p_j^{-2}=(q_j/p_j)M_{jj}^{-1}$ and
$q_j/p_j=1+O(t^\beta)$; both $\beta$ and $.10\delta$ exceed $\kappa$.
Since
\[
 y^\top M^{-1}y=f^\top Mf-\bar m^2\Tr M,
\]
comparison yields
\[
 (F^3)_{ii}\le
 Ct^\kappa f^\top Mf+C\bar m^2\Tr M+Ct\Lambda^+(y_\Sigma).
\]
The row bounds and the active-site count give
\[
 f^\top Mf\le CH_*,\qquad
 \|y\|_2^2\le CH_*,\qquad
 \|y\|_1\le C,\qquad
 \bar m^2\Tr M\le C/N.
\]
For example, the numerator in the last expression is
$(\sum_jq_jF_{ij})^2\le C$, whereas $\Tr M\ge cN$.
Lemma~\ref{lem:fourier-charge} now gives
\[
 (F^3)_{ii}\le C(t^\kappa H_*+tH_*^{1/3}+N^{-1}).
\]
The exponent difference
$\frac13-\frac{2\delta}{3}-\kappa$ is positive, and
$N^{-1}=4^{-20dm}\le C_dt^{1+\delta+\kappa}$ for $n\le m$.
This proves \eqref{eq:sweep-f-cube}.
\end{proof}

The small factor in \eqref{eq:sweep-f-cube} is the extra information
provided by the sharp inverse-form comparison. Bounded row sums alone
would not preserve the fourth-power cap. We next show how these
prestate bounds control an arbitrary family of possible next deletions.

\subsection{The correction to \texorpdfstring{$K$}{K}}

Fix one prestate satisfying Lemma~\ref{lem:sweep-prestate}, and let
$\mathcal B$ be any prestate-measurable subset of its remaining bit ports. For each
$x\in\mathcal B$, all matrices in the following estimates are evaluated
at this same prestate; only its hypothetical fresh scalar $\xi_{0,x}$ varies.
There is no requirement that the hypothetical outcomes be mutually
independent. Every subset and tracked column set below is fixed by the
prestate. All increments are padded by zero at their pivot.

\begin{lemma}[Means and operator squares for $J$]\label{lem:sweep-j}
For every such subset $\mathcal B$,
\begin{equation}\label{eq:sweep-j-mean}
 \left\|\sum_{x\in\mathcal B}\EE J_x\right\|\le Ct^\beta.
\end{equation}
For every collection of allowable hypothetical outcomes and every vector
$u$,
\begin{equation}\label{eq:sweep-j-stack}
 \sum_{x\in\mathcal B}\|J_xu\|^2
 +\sum_{x\in\mathcal B}\|J_x^*u\|^2
 \le Ct^{\delta/2}\|u\|^2.
\end{equation}
In particular $\|J_x\|\le Ct^{\delta/4}$.
\end{lemma}

\begin{proof}
Extend $\mu_x=\EE\gamma_x$ by zero outside $\mathcal B$.
The moment and lag bounds give
$\|\mu\|_\infty\le C\sqrt t\,t^\beta$.
If the omitted pivot rows and columns are restored temporarily, the sum
of the linear means in \eqref{eq:sweep-j-formula} is
$-2h\circ(h\diag(\mu)h)$. Its absolute row sum is at most
\[
 2\|h_i\|_2\|(h\diag(\mu)h)_i\|_2
 \le C\|h\|\|\mu\|_\infty\le Ct^\beta.
\]
Terms incorrectly restored with $x=i$ or $x=j$ have magnitude at most
$CtF_{ij}$ and hence total row and column sums $O(t)$.
For the quadratic mean,
$|\EE\gamma_x^2-a_x^{-1}|\le C\sqrt t$; its absolute row sum is
bounded by
$C\sqrt t\sum_xF_{ix}\sum_jF_{xj}\le C\sqrt t$.
Symmetry and the Schur test prove \eqref{eq:sweep-j-mean}.

For the linear square bound, fix an output $i$ and put $v_j=h_{ij}u_j$.
Restoring the omitted inner index $j=x$ gives
\[
 \sum_x\left|h_{ix}\sum_jh_{xj}v_j\right|^2
 \le T_*^2\sum_jF_{ij}|u_j|^2.
\]
The restored term equals $h_{xx}h_{ix}^2u_x$; its squared sum over
$i,x$ is at most $CtH_*\|u\|^2$.
Deleting the output $i=x$ can only decrease these estimates.
Sum over $i$ and use the bounded column sums of $F$ to obtain
$C(T_*^2+tH_*)\|u\|^2$. The realized scalar $\gamma_x$ is bounded,
so the same estimate applies to the linear part of $J_x$.

For its quadratic part,
\[
 \|l^2(l^2)^\top u\|^2
 \le H_*|(Ku)_x-K_{xx}u_x|^2.
\]
Summation over $x$ gives $CH_*\|u\|^2$ because $\|K\|$ is bounded.
The quadratic coefficient is uniformly bounded. Combining these
estimates proves the first stack in \eqref{eq:sweep-j-stack}.
Every $J_x$ is complex symmetric, so applying the same bound to
$\bar u$ proves the adjoint stack.
\end{proof}

\subsection{The Laplacian increment}

We next control the change in the energy form.  The estimates must hold
when the next pivot is chosen from an arbitrary subset of the remaining
bit ports.  We prove a relative bound for the sum of the conditional
means and two bounds for the sum of the squared output norms.  The latter
combine to give the dual norm needed for the local concentration argument.

\begin{lemma}[Laplacian increment estimates]\label{lem:sweep-laplacian}
Fix a valid prestate as in Lemma~\ref{lem:sweep-prestate}, with current
port set \(V\), and let \(S\) be any subset of its remaining bit ports.
For each \(x\in S\), let \(L_x\) be the Laplacian increment of
Lemma~\ref{lem:single-update}, padded by zero in row and column \(x\).
The conditional expectation for this increment is over its fresh scalar
update, with the prestate and the choice of \(x\) fixed.
Then, with \(\beta=.42\delta\),
\begin{equation}\label{eq:sweep-laplacian-mean}
 \left|\sum_{x\in S}\EE L_x\right|\le Ct^\beta D
\end{equation}
as real quadratic forms on \(V\).  For every common real potential
\(f\in\RR^V\), and for arbitrary choices of the scalar outcomes at the
hypothetical pivots \(x\in S\),
\begin{align}
 \sum_{x\in S}\|L_xf\|_2^2
   &\le Ct^{\delta/2}D(f),&
 \sum_{x\in S}\|L_xf\|_1^2
   &\le Ct^{-1}D(f),                                      \label{eq:sweep-laplacian-stacks}\\
 \sum_{x\in S}W^*(L_xf)
   &\le Ct^{\delta/2}D(f).                                 \label{eq:sweep-laplacian-dual-stack}
\end{align}
The outputs are neutral; in the last display they may equivalently be
padded by zero to the initial ports of the sweep.
\end{lemma}

\begin{proof}
All matrices and coefficients below are evaluated at the fixed prestate.
In particular, the values of \(f\) do not depend on the hypothetical
pivot.  The zero-padded operator \(L_x\) uses \(f|_{V\setminus\{x\}}\)
and has zero output at \(x\).  Whenever we split a difference through
\(x\), we use the value \(f_x\) of this same common potential.

Write \(d_{ij}f=f_i-f_j\).  The uniform bounds on \(p\) give
\begin{equation}\label{eq:sweep-laplacian-unweighted-energy}
 \mathcal E(f):=\sum_{i,j\in V}F_{ij}(d_{ij}f)^2
 \asymp D(f)
 =\frac12\sum_{i,j\in V}p_ip_jF_{ij}(d_{ij}f)^2.
\end{equation}
For brevity put
\[
 c_x=\frac{p_x}{q_x-p_xF_{xx}},\qquad
 \rho_x=|\gamma_x|^2-c_x,\qquad
 \mu_x=\EE\gamma_x .
\]
The prestate bounds and Lemma~\ref{lem:single-update} imply
\begin{equation}\label{eq:sweep-laplacian-coefficients}
 |\gamma_x|+|\rho_x|\le C,\qquad
 |\mu_x|\le C\sqrt t\,t^\beta,\qquad
 |\EE\rho_x|\le Ct^\beta .
\end{equation}
Indeed, \(q_x=p_x+O(t^\beta)\), \(F_{xx}=O(t)\), and hence
\(c_x=1+O(t^\beta+t)\).  The scalar moment bounds give
\(\EE|\gamma_x|^2=1+O(\sqrt t)\) and
\(\mu_x=-\sqrt t(g_{xx}-r_x)+O(t)\).  Since \(\beta<1/2\),
these imply \eqref{eq:sweep-laplacian-coefficients}.

Decompose \(L_x=L_x^{(1)}+L_x^{(2)}\) according to its survivor
edge weights
\begin{equation}\label{eq:sweep-laplacian-edge-split}
 p_ip_j\left[
 -2\Re\!\left(\gamma_x\overline{h_{ij}}h_{ix}h_{xj}\right)
 +\rho_xF_{ix}F_{xj}\right],\qquad i,j\ne x.
\end{equation}
Both parts are real symmetric Laplacians; their edge weights may have
either sign.

\paragraph{The sum of the means.}
Let
\[
 M_\mu=\diag(\mu_x\mathbf1_{\{x\in S\}})_{x\in V},
 \qquad H_\mu=hM_\mu h,\qquad P_p=\diag(p),
\]
and retain the endpoint weights by introducing
\[
 H_p=P_p^{1/2}hP_p^{1/2},\qquad
 A_p=P_p^{1/2}H_\mu P_p^{1/2}
     =H_p\diag(\mu_x\mathbf1_{\{x\in S\}}/p_x)H_p.
\]
First include all intermediate indices \(x\in S\), also when \(x=i\)
or \(x=j\), in the linear mean.  Its quadratic form becomes
\[
 -\Re\sum_{i,j}p_ip_j\overline{h_{ij}}(H_\mu)_{ij}(d_{ij}f)^2.
\]
Writing $M_f=\diag(f)$ and using
$\langle A,B\rangle_{\mathrm{HS}}=\Tr(A^*B)$, this is exactly
\[
 -\Re\langle[M_f,H_p],[M_f,A_p]\rangle_{\mathrm{HS}},
 \qquad
 D(f)=\tfrac12\|[M_f,H_p]\|_{\mathrm{HS}}^2.
\]
The diagonal matrix $M_f$ commutes with
$C_\mu=\diag(\mu_x\mathbf1_{\{x\in S\}}/p_x)$. Therefore
\[
 [M_f,A_p]=[M_f,H_p]C_\mu H_p+H_pC_\mu[M_f,H_p],
\]
and the ideal property of the Hilbert--Schmidt norm gives
\[
 \|[M_f,A_p]\|_{\mathrm{HS}}
 \le2\|H_p\|\|C_\mu\|\|[M_f,H_p]\|_{\mathrm{HS}}.
\]
Since $p$ is bounded above and below,
$\|H_p\|\le Ct^{-1/2}$ and
$\|C_\mu\|\le C\sqrt t\,t^\beta$. Cauchy--Schwarz in the exact
pairing thus bounds its absolute value by $Ct^\beta D(f)$.
In particular this is a signed cancellation with the actual $p_i p_j$
weights retained; no regularity or constancy of $p$ has been used.
For an edge \(ij\), removing the intermediate index \(x=i\), if
\(i\in S\), costs at most
\[
 C|\mu_i h_{ii}|F_{ij}\le Ct^{1+\beta}F_{ij}.
\]
The removal of \(x=j\) has the same bound.  Summing these edgewise
errors costs at most \(Ct^{1+\beta}D(f)\).

The quadratic mean has absolute form bounded by
\[
 Ct^\beta\sum_{x\in S}\sum_{i,j}F_{ix}F_{xj}(d_{ij}f)^2.
\]
Use
\[
 (d_{ij}f)^2\le2(d_{ix}f)^2+2(d_{xj}f)^2
\]
and the bounded row and column sums of the symmetric matrix \(F\).
After extending the \(x\)-sum to \(V\), the last expression is at most
\(Ct^\beta\mathcal E(f)\).  This proves
\eqref{eq:sweep-laplacian-mean} for every subset \(S\).

\paragraph{The linear output stacks.}
For each \(i\in V\), define the complex vector
\[
 u^i_j=p_j\overline{h_{ij}}(f_i-f_j).
\]
The reality of \(f\) allows the real part in
\eqref{eq:sweep-laplacian-edge-split} to be taken after summing over
\(j\).  Therefore, for \(i\ne x\),
\begin{equation}\label{eq:sweep-laplacian-linear-output}
 |(L_x^{(1)}f)_i|
 \le C|h_{ix}|
       \left|\sum_{j\ne x}h_{xj}u^i_j\right|.
\end{equation}
Moreover,
\[
 \sum_i\|u^i\|_2^2\le C\mathcal E(f)\le CD(f).
\]
Temporarily include \(j=x\) in
\eqref{eq:sweep-laplacian-linear-output}.
The bound of Lemma~\ref{lem:sweep-prestate},
\[
 \max_i\|\diag(h_i)h\|^2
 \le C\sqrt{H_*/t}=Ct^{\delta/2},
\]
gives
\begin{equation}\label{eq:sweep-laplacian-linear-two}
 \sum_{x\in S}\sum_i|h_{ix}(hu^i)_x|^2
 \le\sum_i\|\diag(h_i)hu^i\|_2^2
 \le Ct^{\delta/2}D(f).
\end{equation}
For the squared \(\ell^1\) norms, first apply Cauchy--Schwarz in \(i\):
\begin{align*}
 \sum_{x\in S}\left(\sum_i|h_{ix}(hu^i)_x|\right)^2
 &\le
 \sum_x\left(\sum_iF_{ix}\right)\sum_i|(hu^i)_x|^2\\
 &\le C\sum_i\|hu^i\|_2^2
 \le Ct^{-1}D(f).
\end{align*}

It remains to remove the artificially included \(j=x\) terms.
Their absolute values are at most
\[
 C|h_{xx}|F_{ix}|d_{ix}f|
 \le C\sqrt t\,F_{ix}|d_{ix}f|.
\]
Their squared \(\ell^2\) stack is bounded by
\[
 Ct\sum_{i,x}F_{ix}^2(d_{ix}f)^2
 \le Ct\,\max_{i,j}F_{ij}\,\mathcal E(f),
\]
and their squared \(\ell^1\) stack is bounded by
\[
 Ct\sum_x\left(\sum_iF_{ix}|d_{ix}f|\right)^2
 \le Ct\mathcal E(f).
\]
These are absorbed in the preceding bounds, since
\(\max F\le\sqrt{H_*}\).  Dropping the output \(i=x\) only reduces the
majorants.  This proves both required stacks for \(L_x^{(1)}\).

\paragraph{The quadratic output stacks.}
Put \(\varepsilon=\max_{i,j}F_{ij}\le\sqrt{H_*}\).
By \eqref{eq:sweep-laplacian-coefficients}, the absolute output is bounded
by a constant times
\[
 Q_{xi}:=F_{ix}\sum_jF_{xj}|d_{ij}f|.
\]
Here we have included all omitted indices, which only increases the
nonnegative majorant.  Cauchy--Schwarz in \(j\), followed by splitting
differences at \(x\), gives
\begin{align*}
 \sum_{x,i}Q_{xi}^2
 &\le C\sum_{x,i,j}F_{ix}^2F_{xj}(d_{ij}f)^2\\
 &\le C\sum_{x,i}F_{ix}^2(d_{ix}f)^2
   +C\sum_{x,j}\left(\sum_iF_{ix}^2\right)
                       F_{xj}(d_{xj}f)^2\\
 &\le C\varepsilon\mathcal E(f).
\end{align*}
In the last step we used
\(\sum_iF_{ix}^2\le\varepsilon\sum_iF_{ix}\le C\varepsilon\).
For the other stack, the total mass
\(\sum_{i,j}F_{ix}F_{xj}\) is bounded uniformly in \(x\), so
\begin{align*}
 \sum_x\left(\sum_iQ_{xi}\right)^2
 &\le C\sum_{x,i,j}F_{ix}F_{xj}(d_{ij}f)^2\\
 &\le C\mathcal E(f).
\end{align*}
Both estimates remain true with the \(x\)-sum restricted to \(S\).
Since \(\sqrt{H_*}\le t^{\delta/2}\), combining the linear and
quadratic bounds proves \eqref{eq:sweep-laplacian-stacks}.
No factor depending on the number of candidate pivots has been used.

\paragraph{The dual norm stack.}
The Laplacian property implies that every \(y_x=L_xf\) is neutral.
Lemma~\ref{lem:fourier-charge}, including its version for at most two
ports per site, gives
\[
 W^*(y_x)\le\|y_x\|_2^2
      +Cb\|y_x\|_1^{4/3}\|y_x\|_2^{2/3}.
\]
Hölder's inequality over \(x\), followed by
\eqref{eq:sweep-laplacian-stacks} and \(b\le t\), yields
\begin{align*}
 \sum_{x\in S}W^*(L_xf)
 &\le Ct^{\delta/2}D(f)
   +Cb\left(\sum_x\|L_xf\|_1^2\right)^{2/3}
       \left(\sum_x\|L_xf\|_2^2\right)^{1/3}\\
 &\le C\bigl(t^{\delta/2}+t^{1/3+\delta/6}\bigr)D(f)
 \le Ct^{\delta/2}D(f).
\end{align*}
The last inequality uses \(1/3+\delta/6>\delta/2\).
This proves \eqref{eq:sweep-laplacian-dual-stack}.
\end{proof}


\subsection{The transformed diagonal lag}

The next estimate controls increments of the transformed lag
\(\widetilde U=(\Id-K)^{-1}(g_{\rm diag}-r)\). Its cap yields a small
raw diagonal error through Lemma~\ref{lem:sweep-prestate}, preserving the
small conditional means used in the preceding increment bounds. For brevity
write \(u=\widetilde U\) in this subsection.
At a deletion, we freeze the removed coordinate of \(u\), and continue to
update its surviving coordinates.  In particular, the increment is zero
at the deleted coordinate.

Fix the prestate and let \(V\) be its set of ports.  For a prestate-measurable subset
\(S\subseteq V\) of eligible bit ports, an array \(Y=(Y_i^x)\) has row
index \(i\in V\) and hypothetical deletion index \(x\in S\).  Define
\[
 A_j(Y)=\max_{i\in V}
       \left(\sum_{x\in S}|Y_i^x|^j\right)^{1/j},
 \qquad j=1,2.
\]
Every survivor vector associated with deletion \(x\) is extended by zero
at coordinate \(x\).  This convention is used throughout this subsection.

\begin{lemma}[Increments of the transformed lag]\label{lem:sweep-lag}
Let \(t>0\) be sufficiently small, \(\beta=.42\delta\), and
\(H_*=t^{1+\delta}\).  At a fixed prestate, suppose that \(h=h^\top\),
\(g=h/\sqrt t\), \(K_{ij}=h_{ij}^2\), \(F_{ij}=|h_{ij}|^2\), and
\begin{equation}\label{eq:sweep-lag-prestate}
\begin{gathered}
 \max_i\sum_jF_{ij}\le C,\qquad
 \max_i\sum_jF_{ij}^2\le H_*,\qquad
 \|h\|\le Ct^{-1/2},\qquad
 \max_i|h_{ii}|\le C\sqrt t,\\
 T_*:=\max_i\|\diag(h_{ij})_j h\|\le Ct^{\delta/4},\qquad
 \|r\|_\infty\le C,\qquad \|u\|_\infty\le t^\beta.
\end{gathered}
\end{equation}
Here \(u=P(g_{\mathrm{diag}}-r)\). Assume also that
\(P=(\Id-K)^{-1}\) exists and that
\[
 L_P=\max\left(1,\max_i\sum_j|P_{ij}|\right),
 \qquad L_P\le C\log(1/t).
\]
For each \(x\) in any prestate-measurable subset \(S\) of eligible bit
ports, perform the
update of Lemma~\ref{lem:single-update} with a conditionally centered
variable \(\xi_{0,x}\) satisfying
\[
 \EE(\xi_{0,x}^2\mid\mathrm{prestate})=t,
 \qquad |\xi_{0,x}|\le C\sqrt t.
\]
Write \(\Delta u^x\) for the resulting lag increment, padded by zero at
\(x\).  Then, uniformly over \(S\) and over all hypothetical outcomes,
\[
 A_2(\Delta u)\le CL_Pt^{\delta/2}.
\]
Moreover, with all expectations conditional on the prestate,
\[
 \max_i\left|\sum_{x\in S}\EE\Delta u_i^x\right|
 \le C\left[L_P(t^{2\beta}+\sqrt t)+L_P^2t^{\delta/2}\right].
\]
The constants are independent of the number of ports and of \(S\).
\end{lemma}

The row and operator bounds in the statement are supplied by
Lemma~\ref{lem:sweep-prestate}.  The absolute row-sum bound for \(P\) is
proved in Lemma~\ref{lem:inverse-row}; the present argument uses precisely
that displayed bound.  We give the details of the survivor restrictions
and of the nonlinear remainder, since both affect the conditional mean.

\begin{proof}
All matrices without a prime are fixed at the prestate.  Expectations in
the proof are conditional on that prestate.  Set
\[
 R_F=\max_i\sum_jF_{ij},\qquad
 M_F=\max_{i,j}F_{ij}\le\sqrt{H_*}.
\]
Thus \(R_F\le C\).  For a hypothetical deletion \(x\), write
\[
 R_x=V\setminus\{x\},\qquad
 l^x=(h_{ix})_{i\in R_x},\qquad
 k^x=((h_{ix})^2)_{i\in R_x},\qquad a_x=1-K_{xx}.
\]
The diagonal bound gives \(a_x=1+O(t)\), so \(a_x\) is bounded away
from zero.  Lemma~\ref{lem:single-update} gives
\[
 \gamma_x=\frac{\xi_{0,x}+tr_x}
 {\sqrt t\,[1+(\xi_{0,x}+tr_x)g_{xx}]},
\]
\[
 J_x=-2\gamma_x\diag(l^x)h_{R_xR_x}\diag(l^x)
       +(\gamma_x^2-a_x^{-1})k^x(k^x)^\top.
\]

\paragraph{The exact equation.}
Let \(B_x=\Schur_{\{x\}}(\Id-K)\).  In the block decomposition with
deleted coordinate \(x\), the off-diagonal column of \(\Id-K\) is
\(-k^x\).  Eliminating \(U_x\) from \((\Id-K)U=z\) therefore gives
\[
 B_xU_{R_x}=z_{R_x}+k^xz_x/a_x.
\]
On the other hand, the update gives
\[
 \Id-K'=B_x-J_x,\qquad
 z'_{R_x}=z_{R_x}-\gamma_xk^x.
\]
Subtracting the two equations yields
\[
 (B_x-J_x)\Delta U_{R_x}
 =-(\gamma_x+z_x/a_x)k^x+J_xU_{R_x}=:Z^x.
\]
Since \(B_x^{-1}=P_{R_xR_x}\),
\begin{equation}\label{eq:sweep-lag-exact}
 \Delta U_{R_x}
 =(\Id-P_{R_xR_x}J_x)^{-1}P_{R_xR_x}Z^x.
\end{equation}
The inverse in this formula will also follow from the convergent series
below.

\paragraph{Moments and the forcing array.}
The identity \(g_{ii}-r_i=[(\Id-K)u]_i\) and the row-sum bound imply
\(|g_{ii}-r_i|\le Ct^\beta\).  Expanding the denominator of \(\gamma_x\)
using \(|\xi_{0,x}+tr_x|\le C\sqrt t\), centering, and variance \(t\), gives
\[
 \mu_x:=\EE\gamma_x=-z_x+O(t),\qquad
 \EE\gamma_x^2=1+O(\sqrt t).
\]
In particular,
\begin{equation}\label{eq:sweep-lag-moments}
\begin{gathered}
 \max_x|\mu_x|\le C\sqrt t\,t^\beta,\qquad
 \EE\gamma_x+z_x/a_x=O(t),\qquad
 \nu_x:=\EE\gamma_x^2-a_x^{-1}=O(\sqrt t),\\
 |\gamma_x|+|\gamma_x^2-a_x^{-1}|\le C.
\end{gathered}
\end{equation}
For the middle estimate, the extra term is
\(z_x(a_x^{-1}-1)=O(t^{3/2+\beta})\).

Put \(Y^x=Z^x/\sqrt t\).  Its direct forcing satisfies
\[
 A_2\left(-\frac{\gamma_x+z_x/a_x}{\sqrt t}k^x\right)
 \le C\sqrt{H_*/t}=Ct^{\delta/2}.
\]
The linear part of \(J_xu_{R_x}\), apart from its bounded coefficient,
is
\[
 h_{ix}\sum_{j\ne x}h_{ij}h_{xj}u_j.
\]
For fixed \(i\), restore the term \(j=x\) and set
\(D_i=\diag(h_{ij})_j\).  The resulting vector indexed by \(x\) is
\(D_i hD_i u\), of norm at most
\[
 \|D_i h\|\,\|D_i u\|_2\le CT_*t^\beta.
\]
The restored term \(h_{ix}^2h_{xx}u_x\) has \(\ell^2_x\)-norm at most
\(C\sqrt t\sqrt{H_*}\,t^\beta\).  Removing the output \(x=i\)
only reduces this norm.  The quadratic part is bounded by
\[
 \left|h_{ix}^2\sum_{j\ne x}h_{xj}^2u_j\right|
 \le R_Ft^\beta F_{ix},
\]
and hence has \(A_2\)-norm at most \(Ct^\beta\sqrt{H_*}\).
Consequently,
\begin{equation}\label{eq:sweep-lag-forcing-square}
\begin{split}
 A_2(Y)&\le C\left(t^{\delta/2}+t^{\beta+\delta/4}
             +t^\beta\sqrt{tH_*}+t^\beta\sqrt{H_*}\right)\\
       &\le Ct^{\delta/2},
\end{split}
\end{equation}
because \(\beta+\delta/4=.67\delta>.50\delta\).

For the summed mean, the direct forcing costs at most
\(C\sqrt t\sum_xF_{ix}\le C\sqrt t\).  Extend \(\mu_x\) by zero off
\(S\).  Before the survivor omissions, the summed linear part of
\(\EE J_x\) is \(-2h\circ(h\diag(\mu)h)\).  Its absolute row sum is
bounded by
\[
\begin{split}
 \sum_j |h_{ij}|\,|(h\diag(\mu)h)_{ij}|
 &\le \|h_i\|_2\,\|(h\diag(\mu)h)_i\|_2\\
 &\le C\|h\|\,\|\mu\|_\infty\le Ct^\beta.
\end{split}
\]
Contracting against \(u\) contributes another factor \(t^\beta\).
Each omitted set \(x=i\) or \(j=x\) costs at most
\[
 C\|\mu\|_\infty\sqrt t\,R_F\|u\|_\infty,
\]
and their intersection satisfies the same bound.  For the quadratic
part, \(|\nu_x|\le C\sqrt t\), while
\[
 \sum_xF_{ix}\sum_jF_{xj}|u_j|\le R_F^2t^\beta.
\]
Its survivor omissions decrease this absolute bound.  We have proved
\begin{equation}\label{eq:sweep-lag-forcing-mean}
 M_0:=\max_i\left|\sum_{x\in S}\EE Y_i^x\right|
 \le C(t^{2\beta}+\sqrt t).
\end{equation}
This estimate controls the signed sum of the conditional means; it does
not assert an \(A_1\) bound on the forcing itself.

\paragraph{The restricted inverse.}
Let \(M_x=\Id-e_xe_x^*\), and define
\(Q^x=M_xPM_x\).  Since \(Y_x^x=0\), the padded vector
\(P_{R_xR_x}Y^x\) equals \(Q^xY^x=M_xPY^x\).  For an array \(W\),
write \((QW)^x=Q^xW^x\).  Minkowski's inequality gives
\begin{equation}\label{eq:sweep-lag-inverse-arrays}
 A_j(QW)\le L_P A_j(W),\qquad j=1,2.
\end{equation}
For means, the output mask depends on the deletion index.  Precisely,
\[
 \sum_{x\in S}\EE(Q^xY^x)_i
 =\left(P\sum_{x\in S}\EE Y^x\right)_i
   -\boldsymbol{1}_{\{i\in S\}}(P\EE Y^i)_i.
\]
For every hypothetical outcome, \eqref{eq:sweep-lag-forcing-square}
gives \(A_2(Y)\le Ct^{\delta/2}\). Jensen's inequality therefore
bounds the second term by \(CL_Pt^{\delta/2}\). Thus
\begin{equation}\label{eq:sweep-lag-mask}
 \max_i\left|\sum_x\EE(Q^xY^x)_i\right|
 \le L_P M_0+CL_Pt^{\delta/2}.
\end{equation}

\paragraph{Two array bounds for the correction.}
We next control repeated applications of the random matrices \(J_x\).
The bounds below are pathwise and apply to any array \(W\), including
one depending on the same outcomes as \(J_x\).  Write \(A_2(W)=A_W\).
Up to uniformly bounded coefficients, the linear and quadratic outputs
are
\[
 \mathcal L_i^x=h_{ix}\sum_{j\ne x}h_{ij}h_{xj}W_j^x,
 \qquad
 \mathcal Q_i^x=h_{ix}^2\sum_{j\ne x}h_{xj}^2W_j^x.
\]
Both are set to zero when \(i=x\).  Cauchy--Schwarz in \(j\) gives
\[
 \sum_x|\mathcal L_i^x|^2
 \le R_F\sum_{x,j}F_{ix}F_{ij}|W_j^x|^2
 \le R_F^2M_FA_W^2.
\]
Also \(|W_j^x|\le A_W\), so
\[
 |\mathcal Q_i^x|\le R_FA_WF_{ix},\qquad
 \sum_x|\mathcal Q_i^x|^2\le R_F^2H_*A_W^2.
\]
Writing \((JW)^x=J_xW^x\), we obtain
\begin{equation}\label{eq:sweep-lag-J-square}
 A_2(JW)\le C(M_F^{1/2}+H_*^{1/2})A_2(W)
          \le C(H_*^{1/4}+H_*^{1/2})A_2(W).
\end{equation}
All omissions only reduce the nonnegative upper bounds above.

For the \(A_1\) estimate, apply Cauchy--Schwarz over \((x,j)\):
\[
\begin{split}
 \sum_x|\mathcal L_i^x|
 &\le\sum_{x,j}|h_{ix}h_{xj}|\,|h_{ij}W_j^x|\\
 &\le\left(\sum_{x,j}F_{ix}F_{xj}\right)^{1/2}
       \left(\sum_{x,j}F_{ij}|W_j^x|^2\right)^{1/2}
 \le R_F^{3/2}A_W.
\end{split}
\]
The quadratic part satisfies
\[
 \sum_x|\mathcal Q_i^x|
 \le A_W\sum_{x,j}F_{ix}F_{xj}\le R_F^2A_W.
\]
Therefore
\begin{equation}\label{eq:sweep-lag-J-sum}
 A_1(JW)\le CA_2(W).
\end{equation}

\paragraph{The remainder and the final mean.}
Define the columnwise operator \(\mathcal C\) by
\((\mathcal CW)^x=Q^xJ_xW^x\), and put
\[
 q_*=CL_P(H_*^{1/4}+H_*^{1/2}).
\]
Equations~\eqref{eq:sweep-lag-inverse-arrays}
and~\eqref{eq:sweep-lag-J-square} give
\(A_2(\mathcal CW)\le q_*A_2(W)\).  For small \(t\), \(q_*<1/2\).
Equation~\eqref{eq:sweep-lag-exact} thus has the convergent
series solution
\[
 \Delta u=\sum_{k\ge0}\mathcal C^kQY,
 \qquad
 A_2(\Delta u)\le\frac{L_P}{1-q_*}A_2(Y)
                  \le CL_Pt^{\delta/2}.
\]
Apply~\eqref{eq:sweep-lag-J-sum} to the outermost correction in every
term of positive degree.  It gives
\begin{equation}\label{eq:sweep-lag-remainder}
\begin{split}
 A_1(\Delta u-QY)
 &\le CL_P\sum_{k\ge0}A_2(\mathcal C^kQY)\\
 &\le\frac{CL_P^2}{1-q_*}A_2(Y)
 \le CL_P^2t^{\delta/2}.
\end{split}
\end{equation}
Both \(J_x\) and \(Y^x\) depend on \(\xi_{0,x}\); their expectations need
not factor.  The remainder estimate is valid for every collection of
hypothetical outcomes, so averaging it bounds the sum of the absolute
conditional mean remainders.  Combining it with
\eqref{eq:sweep-lag-forcing-mean} and~\eqref{eq:sweep-lag-mask} yields
\[
 \max_i\left|\sum_x\EE\Delta u_i^x\right|
 \le C\left[L_P(t^{2\beta}+\sqrt t)+L_P^2t^{\delta/2}\right],
\]
as required.  No independence between a correction and its forcing was
used.
\end{proof}


\subsection{Fourth powers of a row}

The last estimate applies to two tracked column sets: all remaining
columns, and reserve columns only. In the second case the row itself is
a reserve port. A tracked row is frozen when its port is deleted.
When another port $x$ is deleted, define its proxy increment by taking
the fourth-power change on the surviving tracked columns and ignoring
the favorable loss of the removed column. Thus the accumulated proxy
is an upper bound for the actual change of the tracked row sum.

\begin{lemma}[Fourth-power proxies]\label{lem:sweep-fourth}
At a prestate satisfying Lemma~\ref{lem:sweep-prestate}, fix a surviving
row $i$, a tracked column set $\mathcal C$, and any subset
$\mathcal B$ of possible next bit ports. Let $R_i^x$ be the proxy
increment just defined, with $R_i^i=0$. Then
\begin{equation}\label{eq:sweep-fourth-bounds}
 \left|\sum_{x\in\mathcal B}\EE R_i^x\right|\le Ct^\kappa H_*,
 \qquad
 \sum_{x\in\mathcal B}|R_i^x|^2\le Ct^\kappa H_*^2.
\end{equation}
The second inequality holds for every collection of hypothetical
outcomes. Hence $|R_i^x|\le Ct^{\kappa/2}H_*$.
\end{lemma}

\begin{proof}
The linear part of $R_i^x$, for $x\ne i$, is
\[
 -4\Re(\gamma_x A_{ix}),\qquad
 A_{ix}=h_{ix}\sum_{\substack{j\in\mathcal C\\j\ne x}}
                  F_{ij}\bar h_{ij}h_{xj}.
\]
Cauchy--Schwarz gives
\[
 |A_{ix}|^2\le
 F_{ix}\left(\sum_jF_{ij}^2\right)
          \left(\sum_jF_{ij}F_{xj}\right),
 \qquad
 \sum_x|A_{ix}|^2\le H_*(F^3)_{ii}\le Ct^\kappa H_*^2.
\]

For the mean, extend $\mu_x=\EE\gamma_x$ by zero outside
$\mathcal B\setminus\{i\}$ and put $B=h\diag(\mu)h$.
The Ward row bounds and $\|\mu\|_\infty\le C\sqrt t\,t^\beta$ imply
\[
 \max_j|B_{ij}|\le C\sqrt t\,t^\beta,\qquad
 \|B_i\|_2\le Ct^\beta,\qquad
 \|B_i\|_4\le Ct^{1/4+\beta}.
\]
Therefore the absolute linear mean, summed over $x$, is bounded by
\[
 C\sum_{j\in\mathcal C}|h_{ij}|^3|B_{ij}|+CtH_*
 \le CH_*^{3/4}t^{1/4+\beta}+CtH_*.
\]
The error term accounts for restoring $j=x$: its individual value is
$\mu_xh_{xx}F_{ix}^2$ and its summed absolute value is at most
$CtH_*$.

For complex $a,b$, the elementary fourth-power expansion gives
\[
 \bigl||a-b|^4-|a|^4+4\Re(|a|^2\bar a b)\bigr|
 \le C(|a|^2|b|^2+|b|^4).
\]
Apply this to $a=h_{ij}$ and $b=\gamma_xh_{ix}h_{xj}$.
The sum over $x,j$ of the absolute nonlinear errors is at most
\[
 C\sum_{x,j}(F_{ij}F_{ix}F_{xj}+F_{ix}^2F_{xj}^2)
 \le C((F^3)_{ii}+H_*^2).
\]
After division by $H_*$, the total bias is thus bounded by
\[
 C(t^{\beta-\delta/4}+t+t^\kappa+H_*)\le Ct^\kappa,
\]
since $\beta-\delta/4=.17\delta>\kappa$.
The square stack of the nonlinear errors is bounded by the square of
their absolute sum, and is smaller than $Ct^\kappa H_*^2$.
Combining with the linear square stack proves Lemma~\ref{lem:sweep-fourth}. All bounds
remain valid after restricting either of the two index sets.
\end{proof}

\subsection{The form needed for stopped random ordering}

We collect the estimates in a format that separates algebra from
locality and freshness. Set $\ell=1+\log(1/t)$.

\begin{proposition}[Ordinary-sweep increment estimates]\label{prop:sweep-increments}
Consider a current numerical prefix of a sweep whose initial physical
state satisfies the bounds at a stage boundary in Section~\ref{sec:state},
including self-consistency, and is duplicated as in
Lemma~\ref{lem:duplication}. Assume that this prestate satisfies
\eqref{eq:sweep-global-caps}--\eqref{eq:sweep-row-lag-caps},
the broad bounds on $r,p$, and the active-site count.
For the lag estimates also assume
$\|(\Id-K)^{-1}\|_{\infty\to\infty}\le C\ell$.
Let $\mathcal B$ be any subset of the remaining bit ports and use
fresh scalars satisfying \eqref{eq:sweep-digit-moments}.
Then the estimates of Lemmas~\ref{lem:sweep-j},
\ref{lem:sweep-laplacian}, \ref{lem:sweep-lag}, and
\ref{lem:sweep-fourth} hold before division by the number of remaining
permutation slots.

More precisely, let $\mathcal H$ be a finite-dimensional real Hilbert
space and let $B:\mathcal H\to$ initial-port potentials be a fixed
linear extension map with
\[
 W(Bv)\le C\|v\|_{\mathcal H}^2,\qquad
 \|B^*y\|_{\mathcal H^*}^2\le CW^*(y).
\]
The dual estimate is required for neutral charges $y$.
Interpret $B^*L_xB$ as a self-adjoint operator by the Riesz
identification. The following table gives a bound for the summed
conditional mean, the pathwise square stack, and a single jump:
\[
\begin{array}{c|c|c|c}
 \text{increment}&\text{summed mean}&\text{square stack}&\text{jump}\\ \hline
 J_x&t^\beta&t^{\delta/2}&t^{\delta/4}\\
 B^*L_xB&t^\beta&t^{\delta/2}&t^{\delta/4}\\
 \Delta\widetilde U_i^x&
 \ell(t^{2\beta}+\sqrt t)+\ell^2t^{\delta/2}&
 \ell^2t^\delta&\ell t^{\delta/2}\\
 R_i^x/H_*&t^\kappa&t^\kappa&t^{\kappa/2}
\end{array}
\]
All entries are upper bounds up to a fixed multiplicative constant.
For $J_x$ the square stack includes both orders, as required by
Hermitian dilation.
\end{proposition}

\begin{proof}
Only the pullback and centering assertions remain to explain.
The Schur comparison following \eqref{eq:sweep-global-caps} gives
$D_{\rm cur}(Bv|_{\rm cur})\le CW(Bv)\le C\|v\|_{\mathcal H}^2$.
The form mean and dual square bound of
Lemma~\ref{lem:sweep-laplacian} consequently imply
\[
 \left\|\sum_x\EE(B^*L_xB)\right\|\le Ct^\beta,\qquad
 \sum_x(B^*L_xB)^2\le Ct^{\delta/2}\Id.
\]
The latter follows by applying the assumed dual bound to
$L_x(Bv|_{\rm cur})$ and summing over $x$.
The remaining rows of the table are Lemmas~\ref{lem:sweep-j},
\ref{lem:sweep-lag}, and~\ref{lem:sweep-fourth}.

If there are $M$ remaining slots in a uniformly random ordering, each
conditional mean and second moment is the corresponding sum divided by
$M$; skipped slots contribute zero. Center the joint random
slot-and-scalar increment. For a self-adjoint increment $Y$,
\[
 \EE[(Y-\EE Y)^2]=\EE Y^2-(\EE Y)^2\le\EE Y^2.
\]
The centered jump increases by at most a factor two. Apply this to
the Hermitian dilation for a complex or rectangular increment.
Predictable stopping preserves all these bounds, including the
prospective first breaking step because its prestate still satisfies
the hypotheses.
\end{proof}

These estimates supply the local concentration inputs for an ordinary
sweep. It remains to accumulate them with a failure bound independent of
the torus volume. Section~\ref{sec:local-sweeps} proves the inverse row bound,
constructs fixed local tests, and shows that their budgets imply the
sweep caps. It then applies stopped concentration to those tests and
returns to the physical endpoint described by
\eqref{eq:sweep-reserve-target}.
